%% 06_conclusion.tex — Conclusión y trabajo futuro.

\section{Conclusión y trabajo futuro}
\label{sec:conclusion}

Se ha presentado \herramienta, un motor DAST asíncrono cuya arquitectura separa
explícitamente la planificación de \emph{payloads} y la confirmación de
hallazgos en etapas desacoplables, junto con una metodología de evaluación
reproducible sobre OWASP Benchmark, un oráculo independiente, un diseño RCBD y
un manifiesto con \emph{hash} del corpus. Sobre ese motor se documenta, además,
un agente de triaje asesor especializado mediante una metodología modular de
\emph{fine-tuning} y alineación en tres fases ---saneamiento del conjunto de
datos, adaptación eficiente en parámetros con cuantización, y validación de
seguridad sobre un conjunto disjunto--- reproducible sobre infraestructura de
aceleración estándar con escalado dinámico de parámetros.

El estudio de ablación arroja un resultado matizado. Sobre los
\NumObjetivos{}~objetivos de parámetro \code{GET}, el \code{pipeline} completo
alcanza ya \ValorRecallBudgetDiez{} de exhaustividad con solo
\num{10}~peticiones por parámetro y \emph{satura} en \ValorRecallTecho{} con
\num{20}, con una FPR de \ValorFPRBaseline{} estable en los cuatro
presupuestos evaluados. De las tres optimizaciones evaluadas, solo la
\textbf{priorización por
confianza} produce un cambio con consecuencias prácticas, y acotado a los
presupuestos ajustados: desactivarla cuesta \DeltaRecallNoPriority{} de
exhaustividad y \DeltaPrecisionNoPriority{} de precisión a presupuesto
\num{20}, y nada a partir de \num{50}. El \textbf{entrelazado por contexto} y
la \textbf{confirmación en tiempo de ejecución} no modifican la exhaustividad,
la precisión ni la FPR en este banco; la prueba de Friedman detecta
diferencias en el \emph{orden} de las peticiones
($p = \ValorFriedmanP{}$), pero con tamaños del efecto insignificantes
($\lvert\delta\rvert < \num{0.02}$). La conclusión operativa es que, en un
banco con señales de disparo tan regulares como OWASP Benchmark, el beneficio
de la planificación de \emph{payloads} es de \emph{anticipación} ---mayor área
bajo la curva de detección cuando el presupuesto es escaso--- y no de
\emph{techo} de detección, y que la confirmación en tiempo de ejecución no
penaliza la precisión porque los indicadores crudos ya son fiables en este
escenario.

\noindent\textbf{Trabajo futuro.}
\begin{enumerate}
  \item Repetir el estudio sobre aplicaciones reales y bancos con señales
  ruidosas, donde la confirmación en tiempo de ejecución debería mostrar su
  valor en la reducción de falsos positivos que aquí no fue observable.
  \item Extender el banco a vectores \code{POST}, cabeceras y \emph{cookies}
  para caracterizar categorías hoy infrarrepresentadas.
  \item Ampliar la rejilla a presupuestos menores ($< \num{10}$), donde el
  efecto de la priorización debería ser aún más pronunciado, con réplicas
  para estrechar los intervalos.
  \item Añadir un precalentamiento explícito de la JVM y repetir el análisis de
  latencia absoluta entre configuraciones.
  \item Aprender la política de ordenación \code{freq} de forma adaptativa
  durante el propio escaneo, en lugar de fijarla \emph{a priori}.
  \item Publicar el conjunto completo de manifiestos y trazas JSONL como
  artefacto de investigación verificable.
  \item Cuantificar empíricamente la reducción de falsos positivos del agente
  de triaje frente al motor determinista sobre bancos con señales ruidosas, y
  calibrar la confianza del veredicto para acotar el riesgo de falso negativo.
  \item Extender la metodología de alineación en tres fases a la síntesis de
  \emph{payloads} de confirmación, evaluando si un adaptador especializado
  supera a la lista estática bajo presupuesto ajustado.
\end{enumerate}
